% Chapter 4 -- Debugging Turbo Initech programs (model: TP 6.0 UG Ch 5:
% the taxonomy of bugs, then the tools; Turbo Initech has no integrated
% debugger, ADR-0007 DEC-08, so the tools are the compiler's own listings
% and tracing output).  AN-19.
% Sources (Law 1): os/tps/tps.pas FailLex/FailParse/FailType/GenFail and
% the SKIP lines; os/tps/PARSER-TRACE.md (ENTER/EXIT trace, L = line of the
% lookahead), SYMBOLS-DUMP.md; seed/rt/fileio.asm (FILEIO-ERROR); ADR-0007
% DEC-03 (complete evaluation).  Every diagnostic and listing shown is
% replayed by test-pubs-tps.
\Gates{test-tps-lex-fpc, test-tps-parse-fpc, test-tps-type-fpc, test-tps-parse-mutant, test-compiler, test-pubs-tps}
\chapter{Debugging Turbo Initech programs}
\ix{debugging}
\ix{errors}

A \term{bug} is a difference between what a program does and what you
meant it to do.  Some bugs the compiler finds for you; others only show
when the program runs.  This chapter shows how to find both kinds.

\section{The kinds of errors}
\begin{itemize}
\item A \term{lexical error} is a character or token the compiler cannot
  read: a stray symbol, an unclosed string or comment.
\item A \term{syntax error} is a program that does not follow the grammar:
  a missing semicolon, a missing \rw{end}, a call without parentheses.
\item A \term{semantic error} is a program whose grammar is right but whose
  meaning is not allowed: a name not declared, a value of the wrong type.
\item A \term{run-time error} stops a compiled program: in Turbo Initech,
  only a failed file operation does so.
\item A \term{logical error} produces a wrong result without any message.
\end{itemize}

The compiler stops at the first error of the first three kinds and tells
you which pass found it.

\section{Reading the compiler's messages}
\ix{error messages, reading}
Every message gives the line and the column of the error.  This program
has an apostrophe inside a string:

\Source{err-quote}

\CompilerOutput{err-quote}{ERR}

The string \code{'Didn'} ended at the second apostrophe; the words after it
were read as identifiers, and the apostrophe after \code{you}, in column 32,
began a new string that the end of the line left unclosed.  Write an
apostrophe inside a string as two apostrophes: \code{'Didn''t'}.  The other
passes were skipped.

A syntax error names the token that was found and what the grammar
expected in its place:

\Source{err-noparens}

\CompilerOutput{err-noparens}{ERR}

The compiler read \id{Beep} as the start of an assignment and expected
\code{:=}.  A procedure call needs its parentheses: \code{Beep()}.

A semantic error comes after the syntax pass has printed
\code{TPS-PARSE-OK}:

\Source{err-undeclared}

\CompilerOutput{err-undeclared}{ERR}

The name \id{Totl}, misspelled, is not declared.  The message
\code{semantic rule violation} covers every rule that has no message of its
own; the \doctitle{Programmer's Guide}, Appendix A, lists the rules, and
the column points at the construct concerned.

\section{Seeing how the compiler read a statement}
\ix{parse trace}
The second part of the compiler's output traces every rule of the grammar
it entered and left, with the line it had reached.  When a statement does
not do what you expect, the trace shows how the compiler grouped it.  In
this program, which \rw{if} does the \rw{else} belong to?

\Program{dangle}

\ProgramOutput{dangle}

\CompilerOutput{dangle}{PARSE}

The \textit{else-clause} is entered inside the second \textit{stmt-if}, the
one on line 7, before that \rw{if} is left: the \rw{else} belongs to the
inner \rw{if}, the nearest one.  If you meant it for the outer \rw{if},
put the inner \rw{if} statement between \rw{begin} and \rw{end}.

\section{Checking declarations}
\ix{symbol dump}
The third part lists every name with its kind and type.  For the program
\file{TAX.PAS} of Chapter 2 it is:

\CompilerOutput{tax}{TYPE}

The list shows the parameters of each routine and their types.  A
parameter you meant to be a \rw{var} parameter shows as
\code{value-param}; a function result shows as \code{result}.  Chapter 8
describes every column.

\section{Tracing a running program}
\ix{tracing}
To see what a program does while it runs, write the values of its
variables with \std{writeln}{Writeln}.  A Boolean constant or variable
that switches the tracing on and off saves you from deleting the lines
afterward:

\Program{trace}

\ProgramOutput{trace}

\section{Common mistakes}
\SideHead{Semicolon before else}\code{if A then B; else C} is a syntax
error: the semicolon ends the \rw{if} statement.

\SideHead{Calls without parentheses}Every call has parentheses, even
\code{Beep()}.

\SideHead{Strings first in Write}An argument of \std{write}{Write} that
begins with a string literal is that literal alone; write
\code{Writeln(('A' + S))}, not \code{Writeln('A' + S)}.

\SideHead{Guards with and}\code{(I <= Size) and (A[I] = X)} reads
\code{A[I]} even when \id{I} is beyond \id{Size}.  Use nested \rw{if}
statements.

\SideHead{Overflow}Integer arithmetic wraps around without a message:
$2147483647 + 1$ is $-2147483648$.

\SideHead{Indexes}Indexes are not checked.  An index outside the bounds
reads or overwrites other variables.

\SideHead{Local variables}A local variable has no value until you assign
one; global variables start at zero.

\section{Where to go from here}
Chapter 5 describes how a compiled program becomes a program file, and
what limits apply to it.
